\documentclass[10pt,a4paper,twocolumn]{article}

%-----------------------------
% Paquetes
%-----------------------------
\usepackage[spanish,english]{babel}
\usepackage[utf8]{inputenc}
\usepackage[T1]{fontenc}
\usepackage{newtxtext,newtxmath}
\usepackage[left=1.5cm, right=1.5cm, top=2cm, bottom=2cm, columnsep=0.8cm]{geometry}
\usepackage{titlesec}
\usepackage{graphicx}
\usepackage{hyperref}
\usepackage{fancyhdr}
\usepackage{cite}
\usepackage{parskip}

%-----------------------------
% Formato encabezado
%-----------------------------
\pagestyle{fancy}
\fancyhf{}
\fancyhead[L]{\small IOP Publishing $\vert$ Journal Title}
\fancyhead[R]{\small Journal XX (XXXX) XXXXXX}
\renewcommand{\headrulewidth}{0pt}
\cfoot{\thepage}

%-----------------------------
% Formato títulos
%-----------------------------
\titleformat{\section}
{\large\bfseries}
{\thesection.}{0.5em}{}

\titleformat{\subsection}
{\normalsize\bfseries}
{\thesubsection}{0.5em}{}

\titleformat{\subsubsection}
{\normalsize\itshape}
{\thesubsubsection}{0.5em}{}

%-----------------------------
% Documento
%-----------------------------
\begin{document}

\selectlanguage{english}

% El entorno @twocolumnfalse permite poner contenido a una sola columna (el título y abstract)
% y que el resto del documento fluya en las dos columnas puestas en el documentclass.
\twocolumn[
  \begin{@twocolumnfalse}
    
    % Logos u opcional
    \noindent {\small \textbf{IOP Publishing} \hfill \textbf{Journal Title}} \\
    \noindent {\small Journal XX (XXXX) XXXXXX \hfill \url{https://doi.org/XXXX/XXXX}}
    
    \vspace{1.5cm}
    
    % Título
    {\Huge \textbf{Diseño del  montaje experimental para la reproducción y detección del fenómeno de la sonoluminiscencia}}

    
    \vspace{0.8cm}
    
    % Autores
    {\large \textbf{Juan José Montoya Sánchez$^1$, Jose David Ortiz Campo$^1$}}
    
    \vspace{0.3cm}
    
    {\small $^1$ Instituto de Física, Facultad de Ciencias Exactas y Naturales, Universidad de Antioquia, Medellín, Colombia}
    
    \vspace{0.4cm}
    
    % Correspondencia y fechas
    {\small \textbf{E-mail:} \url{juan.montoya110@udea.edu.co}, \url{jose.ortizc@udea.edu.co}}
    
    \vspace{0.8cm}
    
    % Abstract
    \textbf{\large Abstract} \\
    Sample text inserted for illustration. Replace with abstract text. Your abstract should give
    readers a brief summary of your article. It should concisely describe the contents of your
    article, and asaaaaaaaaaaaaaaaaaaaaaaa
    
    \vspace{0.3cm}
    
    \textbf{Keywords:} term, term, term
    
    \vspace{1cm}
  \end{@twocolumnfalse}
]

%-----------------------------
% Secciones a doble columna
%-----------------------------
\section{Introduccion}

El fenómeno de la sonoluminiscencia consiste en la transducción de energía acústica en emisiones de luz, producida por el colapso violento y periódico de burbujas de gas atrapadas en el interior de un líquido sometido a ondas ultrasónicas \cite{gaitan1992, brenner2002}. Aunque a nivel físico general este proceso se caracteriza por generar condiciones termodinámicas extremas, desde una perspectiva metrológica sus rasgos más relevantes radican en su naturaleza transitoria extrema. Los destellos luminosos emitidos poseen duraciones ultracortas, frecuentemente registradas por debajo de los 100 picosegundos, y exhiben una intensidad sumamente baja, llegando a emitir apenas del orden de un millón de fotones por pulso \cite{barber1992, gompf1997}. La caracterización de eventos que combinan de manera simultánea una alta velocidad con una baja intensidad lumínica impone un desafío instrumental dado que los sensores fotónicos enfrentan un límite físico fundamental entre la resolución temporal y la relación señal-ruido. Por ejemplo, equipos altamente especializados como las cámaras de barrido (streak cameras) logran resoluciones temporales en el rango de los picosegundos, pero su baja eficiencia cuántica ante haces de luz tan tenues obliga frecuentemente a promediar la emisión de miles de pulsos sucesivos para obtener una medición viable, perdiendo así la información dinámica de cada evento individual \cite{pecha2000}. De igual forma, las cámaras de alta velocidad convencionales, basadas en sensores CCD o CMOS, resultan útiles para registrar la dinámica macroscópica y la oscilación radial de la burbuja, pero carecen de la velocidad de obturación y la sensibilidad requerida (conteo de fotón único) para medir la curva de intensidad del destello aislado sin que la señal quede completamente enmascarada por el ruido térmico y de lectura del propio sensor.

Para superar las restricciones inherentes a la captura de eventos ópticos ultrarrápidos y tenues, el estado del arte instrumental se ha focalizado en el desarrollo de detectores de estado sólido de fotón único y arquitecturas de registro temporal de extrema precisión \cite{becker2005}. Históricamente, la técnica de Conteo de Fotones Individuales Correlacionados en el Tiempo (TCSPC) se ha consolidado como el estándar metrológico para alcanzar resoluciones en el régimen de los picosegundos; sin embargo, su aplicación clásica suele requerir la integración estadística de múltiples ciclos de excitación, limitando la caracterización de eventos transitorios únicos e irrepetibles \cite{bruschini2019}. En la actualidad, este nicho tecnológico está liderado por la adopción de matrices de Diodos de Avalancha de Fotón Único (SPAD arrays) y Fotomultiplicadores de Silicio (SiPM), dispositivos capaces de operar en modo Geiger para registrar la llegada de fotones individuales con un nivel de ruido de lectura virtualmente nulo y tasas de muestreo en el rango de los gigahercios \cite{becker2005, mccarthy2013}. Para garantizar la fidelidad de las mediciones, resulta indispensable someter estos equipos a procesos de calibración rigurosos mediante sistemas capaces de generar pulsos lumínicos artificiales altamente controlados, los cuales deben emular con exactitud las exigentes condiciones de velocidad y atenuación que presentan los fenómenos físicos naturales de esta índole.

En este contexto, el presente proyecto se centra en el diseño, caracterización y validación de un montaje experimental optoelectrónico concebido para la detección y medición de destellos lumínicos con características análogas a la sonoluminiscencia. Como paso previo a la implementación de sensores de alta complejidad, se propone el desarrollo de un entorno de simulación basado en hardware accesible, utilizando un microcontrolador ESP32 y un diodo emisor de luz (LED) confinados en una cámara de aislamiento oscuro para generar pulsos ultrarrápidos y de intensidad modulable. Para la adquisición de los eventos, se integra una cámara web comercial, lo que impone un desafío metrológico de submuestreo temporal severo frente a la latencia de microsegundos del pulso luminoso generado. Para sortear esta limitación física del sensor de adquisición convencional, se introduce y evalúa una metodología algorítmica basada en la superposición de fotogramas. Mediante el registro asíncrono y la integración estadística de múltiples eventos periódicos, este método permite reconstruir de forma indirecta la curva transitoria de emisión del destello. Así, este trabajo busca validar la eficacia de técnicas computacionales de procesamiento de imágenes acopladas a instrumentación estándar, estableciendo las bases tecnológicas y metodológicas fundamentales para la futura medición experimental de la sonoluminiscencia real.

\section{Metodología}

El diseño y validación del sistema de detección para eventos luminosos transitorios se estructuró en tres fases experimentales correlacionadas. Dada la brevedad del estímulo luminoso frente a las limitaciones impuestas por el submuestreo de un hardware de adquisición de video convencional, la metodología descarta una grabación directa de alta velocidad y emplea, en su lugar, un modelo estadístico de coincidencias cruzadas. Inicialmente, se caracterizó el nivel de respuesta y la latencia de un canal optoelectrónico de estado sólido auxiliar. Seguidamente, se determinaron los umbrales de sensibilidad del sensor fotográfico principal dentro de un cerramiento aislado. Finalmente, se diseñó e implementó un algoritmo de validación dual encargado de reconstruir la dinámica transitoria del destello mediante técnicas de promediado de fases apoyadas en un emparejamiento temporal asíncrono.

\subsection{Caracterización temporal del sistema analógico emisor-sensor}
Para establecer una referencia de sincronización firme (ancla temporal) que operase a mayor velocidad que la tasa de fotogramas del sensor de video, se implementó un circuito guiado por un microcontrolador ESP32. Este dispositivo lógico se encargó de emitir señales moduladas por ancho de pulso (PWM) hacia un diodo LED, y de capturar simultáneamente la refracción de luz usando un fototransistor analógico OP598 acoplado a un convertidor analógico a digital (ADC). 

Para cuantificar los límites empíricos de adquisición de tal hardware, se diseñó un protocolo de sondeo continuo (\textit{polling}). Este procedimiento evaluó la diferencia temporal mínima entre la orden dictada para la excitación del LED y la confirmación de lectura en el fototransistor tras superar su propio umbral basal. Al determinar la latencia real y la resolución efectiva de adquisición de este circuito de estado sólido, se comprobó su viabilidad metrológica para actuar como árbitro temporal. En consecuencia, el registro instantáneo obtenido por el fototransistor fue establecido formalmente como el ancla de sincronización para los pasos subsecuentes del estudio.

\subsection{Acondicionamiento optoelectrónico y detectabilidad visual}
Una vez validada la referencia de hardware rápida, se evaluaron las cotas de detectabilidad espacial del sistema comercial de captura de video (cámara web). El ensayo tuvo lugar en el interior de una cámara opaca acondicionada visualmente frente al ruido ambiental, escenario necesario para lograr que eventos constituidos por ínfimas tasas fotométricas fuesen disociables frente a los falsos positivos producidos por la excitación o el ruido térmico del propio sensor CMOS de la cámara.

El procedimiento computacional consistió en analizar consecutivamente el espacio ininterrumpido. Se cuantificó la base de iluminación ambiental latente tomando como base métrica el promedio del ruido provocado en estado de reposo (\textit{dark-control}). Posteriormente, se clasificó como ``fotograma detectado'' a aquel donde la intensidad de los píxeles integrados sobrepasara un porcentaje diferencial mínimo sobre dicho nivel estático basal. Tras el acondicionamiento, se corrieron barridos automatizados (\textit{sweeps}) aislados que iteraban las variables configurables del ensayo, a saber, el ciclo de trabajo o esfuerzo de ignición del LED, la duración neta del pulso emitido, y los factores de exposición fotográfica de la cámara. Este proceso permitió establecer sistemáticamente de qué magnitudes requeriría un destello puntual para sembrar huella analizable en el cauce temporal del flujo de video.

\subsection{Reconstrucción estadística asíncrona por superposición espacial}
La última fase experimental unificó ambos sistemas para plantear una solución algorítmica al conflicto del submuestreo temporal. Al tratarse de un sensor de video que posee ventanas de inactividad mientras procesa o renueva un cuadro integrado, existía la probabilidad inherente de que sucesos rápidos transcurrieran dentro de la fracción pasiva del obturador ciego. Para disipar el sesgo y anular un posible emparejamiento estático inercial (\textit{phase lock}) a los refrescos constantes del entorno, se efectuaron lanzamientos duales asíncronos bajo un esquema estocástico con un tren de múltiples igniciones separadas por pausas de tiempos virtual y matemáticamente aleatorios.

Durante estas sesiones, el fototransistor analógico recopiló la hora y latitud precisos en que su componente físico asimiló el resplandor principal. Concluida la captura paralela en ambas vías, un análisis computacional externo (\textit{offline}) se encargó del rastreo. Se revisó de manera algorítmica la línea de tiempo del código de video, determinando la ventana probabilística más cercana al marcador temporal del OP598. Este análisis emparejó los datos de la forma más idónea posible para definir cuál en todo el lote era el instante de mayor fulgor (fotograma \textit{hit}). Con este cruce ejecutado y la confirmación en el punto temporal correcto, la matriz extrajo de manera simétrica los fotogramas circunscritos antes y después del mismo. Finalmente, al apilar por capas sucesivas todos los micro-desgloses recopilados dentro del grueso estocástico y promediar estadísticamente el conjunto píxel por píxel, fue posible consolidar imágenes continuadas de la respuesta estructural fenoménica, posibilitando una aproximación cualitativa al evento real por medio de instrumentación convencional subóptima.
%-----------------------------
% Acknowledgements
%-----------------------------
\section*{Acknowledgements}

Proin pharetra nonummy pede. Mauris et orci.

%-----------------------------
% Referencias
%-----------------------------
\begin{thebibliography}{99}

\bibitem{gaitan1992}
Gaitan D F, Crum L A, Church C C and Roy R A 1992
\textit{Sonoluminescence and bubble dynamics for a single, stable, cavitation bubble.} The Journal of the Acoustical Society of America, \textbf{91}(6), 3166-3183.

\bibitem{brenner2002}
Brenner M P, Hilgenfeldt S and Lohse D 2002
\textit{Single-bubble sonoluminescence.} Reviews of Modern Physics, \textbf{74}(2), 425-484.

\bibitem{barber1992}
Barber B P, Hiller R, Lofstedt R, Putterman S J and Weninger K R 1992
\textit{Resolving the picosecond characteristics of synchronous sonoluminescence.} Physical Review Letters, \textbf{69}(26), 3839-3842.

\bibitem{gompf1997}
Gompf B, Günther R, Nick G, Pecha R and Eisenmenger W 1997
\textit{Resolving sonoluminescence pulse width with time-correlated single photon counting.} Physical Review Letters, \textbf{79}(7), 1405-1408.

\bibitem{pecha2000}
Pecha R and Gompf B 2000
\textit{Microdynamics of sonoluminescence.} Physical Review Letters, \textbf{84}(6), 1328-1330.

\bibitem{becker2005}
Becker W 2005
\textit{Advanced Time-Correlated Single Photon Counting Techniques.} Springer Series in Chemical Physics, Vol. 81. Springer-Verlag.

\bibitem{bruschini2019}
Bruschini C, Homulle H, Antolovic I M, Burri S and Charbon E 2019
\textit{Single-photon avalanche diode imagers in biophotonics: review and outlook.} Light: Science \& Applications, \textbf{8}(1), 87.

\bibitem{mccarthy2013}
McCarthy A, Ren X, Della Frera A, et al. 2013
\textit{Kilometer-range depth imaging at 1550 nm wavelength using an InGaAs/InP single-photon avalanche diode detector.} Optics Express, \textbf{21}(19), 22098-22113.

\end{thebibliography}

\end{document}
